\begin{theorem}[Real-time work law]\label{sw:thm:law}
Let $\Phi$ be time-covariant. In real time, the optimal work per use at vanishing error is $\kB T\ln 2\,\Trt(\Phi)$. Moreover, for every fixed $\delta>0$, any implementation whose controller starts in a thermal state and that draws at most $n\kB T\ln 2\,(\Trt(\Phi)-\delta)$ of deterministic work, meaning that its battery ends in a definite energy level with probability tending to one, has error tending to one as $n\to\infty$.
\end{theorem}

At zero Hamiltonian the first statement, in purity units, is in our earlier work~\cite{Memory,ImmediateDelivery,Douglas26closed}, and so is the second at zero rate~\cite{Flagship}. In batch, the converse of Faist, Berta and Brand\~ao already holds at every fixed error below $1/2$~\cite[Prop.~4.2]{FBB21}.

The rate is a maximum over inputs of a minimum over implementations, while a device is built once and must serve every input and every observer. Part of Theorem~\ref{sw:thm:law} is that the two orders agree: one device costs no more per use than the hardest input, as Faist, Berta and Brand\~ao found in batch, for every channel with Gibbs-preserving maps and for time-covariant channels with thermal operations~\cite{FBB19,FBB21}. For this real-time functional the interchange holds for every channel (Lemma~\ref{sw:lem:minimax}); time covariance lets the construction of Theorem~\ref{sw:thm:law} attain it with a controller that returns exactly, and Theorem~\ref{sw:thm:clock} attains it for every channel with one that returns approximately. In batch with a returned clock, the analogous interchange between choosing a device and choosing an input can fail. For example, $X$-dephasing on a qubit with energy gap $\Delta$ costs $\Delta/2$ per use for any device that is correct on every input (Appendix~\ref{app:batch}), while each known input can be served at no more than the thermodynamic capacity, by the fixed-input construction of Faist, Berta and Brand\~ao~\cite[Prop.~8.3]{FBB21}, whose coherence sources are sublinear in size; they treat them as disposable, but the sources in fact return approximately (Appendix~\ref{app:batch}).

Input by input, the two costs differ in a single correlation term. For a classical channel it vanishes on inputs diagonal in the channel's basis, and those inputs suffice for both optimized rates; on coherent inputs it can be positive (Appendix~\ref{sw:app:classical}). For a partial quantum reset the optimized rates differ.

\begin{theorem}[Classical channels pay no surcharge; partial reset does]\label{sw:thm:classical}\leavevmode
\begin{enumerate}[label=(\alph*), leftmargin=*]
\item If $\Phi$ is classical, meaning that it measures in an eigenbasis of $H$ and prepares states diagonal in that basis, then $\Trt(\Phi)=\Tcap(\Phi)$ at every temperature.
\item Let $\Phi$ be amplitude damping with probability $p$ on a qubit with energy gap $\Delta$. Then $\Trt(\Phi)=\max_{0\le y\le p}\bigl[h(y)-\beta\Delta y/\ln 2\bigr]$, where $h$ is the binary entropy. For $0<p<1$, $\Trt(\Phi)>\Tcap(\Phi)$ at every temperature. For $p=1$, $\Trt(\Phi)=\Tcap(\Phi)$.
\end{enumerate}
\end{theorem}

Amplitude damping with $p=1$ is a full reset to the ground state, the process of Landauer's principle, and timing does not change its cost. Partial reset is different. Its output stays correlated with what the process leaves behind, and in real time that correlation costs work at every temperature (Figure~\ref{sw:fig:damping}).

The surcharge can be as large as the whole cost. For exact implementations we count the free energy $w$ drawn from the battery and the controller (Section~\ref{sw:sec:law}); with a mechanical battery whose entropy changes by $o(n)$ and a controller that returns, this is the work up to $o(n)$.

\begin{theorem}[Werner--Holevo]\label{sw:thm:werner}
Let $W(\rho)=(\Tr(\rho)\,I-\rho^{T})/2$ be the Werner--Holevo channel on a qutrit with zero Hamiltonian. The free energy $w$, in bits, needed for $n$ uses satisfies
\begin{center}
\begin{tabular}{lcc}
\toprule
 & exact & error $n^{-2}$\\
\midrule
real time & $w\ge n\log_23$ & $w\ge n\log_23-O(\log n)$\\
batch & $w\ge n\log_23$ & $w=o(n)$ achievable\\
\bottomrule
\end{tabular}
\end{center}
Both exact bounds are attained with an information battery, and the real-time work rate is $\Trt(W)=\log_23$.
\end{theorem}

In purity units these statements are in our earlier work~\cite{ImmediateDelivery}, and the batch rate of zero follows from the theorem of Faist, Berta and Brand\~ao; the statements here are about free energy. The batch saving comes from tolerating a vanishing error, not from an exact refund.

\paragraph{Further results.} Four further results complete the picture. At zero error the saving of waiting disappears for every extreme channel, such as Werner--Holevo and amplitude damping with $0<p<1$: a batch device then pays the full real-time rate (Proposition~\ref{sw:prop:extreme}, Section~\ref{sw:sec:werner}). In real time a channel costs nothing exactly when it is a limit of thermal operations (Theorem~\ref{sw:thm:free}, Section~\ref{sw:sec:free}); this is a robust, work-cost form of the criterion of Lie et al.~\cite{LieEtAl26}, and it identifies the processes that are free in batch but cost work in real time. A device whose clock must return exactly, and whose battery holds no phase reference, cannot run a channel that is not time-covariant, whatever work is spent (Proposition~\ref{sw:prop:clock}), while a clock that may return approximately runs every channel at the rate $\Trt$ (Theorem~\ref{sw:thm:clock}, Section~\ref{sw:sec:free}). Finally, whether a generalized dephasing channel is free in real time is undecidable in a fixed dimension (Theorem~\ref{sw:thm:undecidable}, Section~\ref{sw:sec:undecidable}).

\paragraph{Relation to prior work.} The batch cost is the thermodynamic capacity of Faist, Berta and Brand\~ao~\cite{FBB19,FBB21}, which had appeared as the work value of channel uses~\cite{NGP15}. Their converse rests on single-shot work costs for a fixed input: the smooth max-entropy of the discarded information conditioned on the output, for degenerate Hamiltonians~\cite{FDOR15}, and the coherent relative entropy in general~\cite{FR18}. The last step of their thermal-operation construction, erasing the environments with the outputs as side information, a universal form of the conditional erasure of~\cite{delRio11} built by position-based decoding~\cite[Sec.~7]{FBB21}, is the step a real-time device cannot take. They suggested replacing the diamond norm by distinguishability measures for processes with memory~\cite{Chiribella08}, noting that this ``would be particularly relevant for settings with memory effects, for instance implementations of gates in a quantum computer''~\cite{FBB19}. Our real-time error is a measure of this kind~\cite{Chiribella08,GutoskiWatrous07}.

Ryb\'ar and Ziman introduced repeatable channels, implemented again and again by one device with a finite memory and no resets, so that every use acts as the same channel. They showed that every mixture of unitaries is repeatable and that no non-unital channel is, because, when the same input $\rho$ is fed repeatedly, the memory's entropy must grow by at least $S(\rho)-S(\Phi(\rho))$ per use, which for a non-unital channel is positive at the maximally mixed input~\cite{RybarZiman}; Lemma~\ref{sw:lem:accumulation} is the conditional form of their entropy count. They showed that for qubits repeatability is equivalent to unitality, and left the general unital case open. The Werner--Holevo channel shows that unitality is not enough in dimension three. Lemma~\ref{sw:lem:accumulation} uses only the output of each use on a fresh input, so it applies to their definition: by Lemma~\ref{sw:lem:invertible}, each exact use on half of a maximally entangled pair adds $\log_23$ bits to the memory (proof of Theorem~\ref{sw:thm:werner}). A memory without resets therefore needs dimension at least $3^n$ for $n$ uses, which their $n$-copy construction attains, and no finite memory repeats the channel indefinitely. The question they close with, whether a unital channel that is not a mixture of unitaries can be repeatable, has a positive answer: a Schur channel with a catalytic dilation~\cite[Prop.~15]{LieEtAl26}, such as the six-level example of Haagerup and Musat~\cite[Ex.~3.3]{HaagerupMusat11}, is repeatable with one finite memory, but its outputs can then be correlated, and such a device need not meet our real-time error (Section~\ref{sec:related}). Mausberg asks what a closed machine needs to transform fresh inputs repeatedly while its memory returns after every use~\cite{Mausberg26}. Exact return rules out purifying maximally mixed inputs, and exact purification of $n$ qubits with return tolerance $\delta$ needs $n/\delta+O(n)$ memory qubits. In the other direction, a memory can return exactly while its outputs, each correct, are perfectly correlated, as in the random-unitary devices of Ryb\'ar and Ziman, whereas independent outputs need $\lceil n/2\rceil$ memory qubits. Our real-time error asks for the joint output, against every observer. Work extraction under causal constraints has been studied for multitime processes~\cite{ZambonAdesso}, for temporally correlated sequences of states~\cite{HuangTOFE} and for agents that act on environments with memory~\cite{Fiderer26}; the online execution of classical strategies with memory has a work cost of its own, which quantum memory can reduce below that of every classical agent whose memory is a function of the past~\cite[Result~2]{Thompson25}. Here we price the implementation of a memoryless channel with quantum inputs and outputs instead.

Quantum processes are known to carry costs that classical ones can avoid: in operations that must act correctly on each of several signal states, not only on their average, which in general generate heat above the bound set by the average state, although classical operations, resets and unitaries do not~\cite{BedinghamMaroney}; in resetting almost every state~\cite{RiechersGu21}; and in the q-machine simulators of a classical process, which can dissipate less than its $\varepsilon$-machine but, like their time reverses, cannot both compress its memory and dissipate nothing, a restriction not established for every quantum realization of the process~\cite{LoomisCrutchfield20,LoomisCrutchfield20PRL}; here the excess comes from timing alone. Classical processes can incur online costs under other constraints, such as a fixed protocol, a modular design or periodic operation~\cite{KolchinskyWolpert21,BMC18,OuldridgeWolpert23}. Our devices may carry any memory and act jointly on everything they hold, so these constraints are absent; the real-time surcharge is instead a modularity cost in time, since the device can no longer act on the outputs it has released (Section~\ref{sw:sec:discussion}).

Our earlier work proves the zero-Hamiltonian, purity-unit forms of the rate statement in Theorem~\ref{sw:thm:law}, of Theorem~\ref{sw:thm:werner} and of Theorem~\ref{sw:thm:free}~\cite{Memory,ImmediateDelivery,Douglas26closed}, a strong converse at zero rate, and undecidability for dephasing channels of unbounded dimension~\cite{Flagship}. The finite-temperature law with work in place of purity, the converse at every fixed error and every rate below $\Trt$, Theorem~\ref{sw:thm:classical}, the robust finite-temperature form of Theorem~\ref{sw:thm:free}, Theorem~\ref{sw:thm:clock}, and the fixed-dimension and non-enumerability statements of Theorem~\ref{sw:thm:undecidable} first appeared in~\cite{SaveWork}.

\paragraph{Outline.} Section~\ref{sw:sec:setting} sets up the model. Section~\ref{sw:sec:law} proves Theorem~\ref{sw:thm:law}. Section~\ref{sw:sec:classical} treats classical channels and reset, Section~\ref{sw:sec:werner} the Werner--Holevo channel, Section~\ref{sw:sec:free} what is free and what a clock can do, and Section~\ref{sw:sec:undecidable} undecidability. Section~\ref{sw:sec:discussion} discusses an experimental test and open problems. The longer constructions are in the appendices. Section~\ref{sec:law} places this law among general schedules of arrivals, releases and loans.


